\title{Same Loss and Gradient Covariance, Different Stationary Errors in Adaptive Optimization}
\titlerunning{Pairing and Stationary Error in Adaptive Optimization}
\author{Ihor Kendiukhov}
\authorrunning{Ihor Kendiukhov}
\institute{Ihor Kendiukhov (corresponding author) \at
Institute of Medical Genetics and Applied Genomics, University of T\"ubingen, T\"ubingen, Germany\\
\email{ihor.kendiukhov@student.uni-tuebingen.de}}
\date{}
\hypersetup{pdftitle={Same Loss and Gradient Covariance, Different Stationary Errors in Adaptive Optimization},pdfauthor={Ihor Kendiukhov}}
\maketitle
\begin{abstract}
The average loss and gradient covariance need not determine the long-run
error of an adaptive gradient method. We prove this for AdamW using two
quadratic problems with two interacting parameters and smooth sampling
distributions. Every sampled loss is strongly convex. The two problems
have the same average loss, the same separate distributions of sample
curvature and gradient noise at the average-loss minimum, and the same
gradient covariance at every pair of parameter values. Only the pairing
of curvature and noise differs.

With identical fixed optimizer settings, stationary error decreases in
proportion to the learning rate in one problem but remains bounded away
from zero in the other. Plain stochastic gradient descent has exactly
the same stationary error on both problems. An exact formula explains
how AdamW weights curvature according to the noise magnitude observed
with it. We identify conditions restoring vanishing error, give the
sharp error range in one dimension, extend the construction to a range
of pairings, and quantify a sequential long-memory limit. Simulations of
a discrete example support the mechanism.

The result provides a controlled test of theories of adaptive optimization:
loss and covariance summaries alone can miss a change in the order of
stationary error. The formula identifies additional dependence that such
theories need to retain.
\end{abstract}
\keywords{AdamW \and Stochastic approximation \and Invariant distributions \and Stochastic optimization \and Curvature--noise dependence}
\subclass{62L20 \and 90C15 \and 60J05}
